\documentclass{article}
% Shared layout. Keep benchmark prose and result-specific plots outside this folder.
\PassOptionsToPackage{numbers}{natbib}
\usepackage[preprint]{neurips_2026}
\usepackage[T1]{fontenc}
\usepackage[hidelinks]{hyperref}
\usepackage{booktabs,amsmath,microtype,pgfplots}
\pgfplotsset{compat=1.17}
\usetikzlibrary{patterns}
\definecolor{okblue}{HTML}{0072B2}
\definecolor{okorange}{HTML}{D55E00}
\pgfplotsset{every axis/.append style={font=\small, tick label style={font=\small}, width=\linewidth, height=5cm, grid=major}}

\usepackage{graphicx}
\newcommand{\PaperTitle}{Memory Conflict: Auditing Household Decisions Under Changing Preferences}
\newcommand{\PaperAbstract}{Stored preferences can become invalid without disappearing. Memory Conflict evaluates whether household assistants act on current preferences while preserving privacy, booking state, and authority boundaries. We audit 40 synthetic episodes across 6 conflict types with 12 scripted conditions and 960 runner episodes including repeats. Reference actions achieve pass rate 1.000 [0.912, 1.000]; empty output achieves 0.000 [0.000, 0.088]. Moving clarification after booking also achieves 1.000 [0.912, 1.000], exposing a temporal limit of final-state scoring. These are 95\% Wilson intervals over tasks, not population guarantees. A clearly simulated population sweep scores 30000 traces under injected errors and perfect control repairs. The sweep estimates sensitivity under explicit assumptions, not model performance. No learned model or human participant is evaluated.}
\newcommand{\AuthorOrcid}{https://orcid.org/0009-0005-9964-6841}
\newcommand{\TaskCount}{40}
\newcommand{\TrapCount}{6}
\newcommand{\ConditionCount}{12}
\newcommand{\EpisodeCount}{960}
\newcommand{\SimCount}{30000}
\newcommand{\SeedCount}{30}
\newcommand{\Good}{1.000 [0.912, 1.000]}
\newcommand{\Empty}{0.000 [0.000, 0.088]}
\newcommand{\NoAsk}{0.825 [0.681, 0.913]}
\newcommand{\NoCancel}{0.825 [0.681, 0.913]}
\newcommand{\Stale}{0.300 [0.181, 0.454]}
\newcommand{\Leak}{0.875 [0.739, 0.945]}
\newcommand{\Late}{1.000 [0.912, 1.000]}
\newcommand{\AskFail}{0.000 [0.000, 0.354]}
\newcommand{\AskPass}{1.000 [0.646, 1.000]}
\newcommand{\SimNone}{0.478 [0.455, 0.502]}
\newcommand{\SimAll}{1.000 [1.000, 1.000]}
\newcommand{\SimGain}{0.522 [0.498, 0.545]}

\title{\PaperTitle}
\author{Pooria Arab\\Independent researcher\\\href{\AuthorOrcid}{\AuthorOrcid}}
\begin{document}
\maketitle
\begin{abstract}\PaperAbstract\end{abstract}
\section{Why this matters}
A household assistant can remember a statement correctly and still make the wrong decision. A preference can expire, another member can have authority, and a private fact can remain useful without becoming shareable. Retrieval quality alone does not settle these questions. Once an assistant books a service, the relevant outcome includes the action and its effects on other household members.\par
Consider an assistant that receives a new restaurant preference while an old reservation remains active. Selecting the right restaurant does not remove the old reservation. An assistant can therefore give a correct explanation while leaving an incorrect state. A second assistant might choose the same restaurant after revealing a private plan to another member. Both failures require evidence beyond the final sentence.\par
Memory Conflict defines a small setting for this distinction. Its target construct is \emph{acting on valid household information while preserving the specified booking, privacy, and authority constraints}. Validity here comes from explicit task rules. It does not mean that the benchmark resolves real family disputes. The rules identify an expected option, required cancellations, restricted disclosures, and situations that require a question.\par
This paper contributes an audit of that construct, a set of controlled failures, and a clearly simulated study of perfect repairs. The audit asks whether the score responds to the changes it should detect. It also asks which changes the score misses. The simulated study shows how task coverage limits the apparent value of a targeted control. Neither study measures a trained agent's memory ability.\par
The central distinction is between an intended rule and its observed test. The tasks describe confirmation before spending. The current score checks whether a qualified person was asked somewhere in the trace. Those conditions differ. Our timing control makes that difference visible without changing the booking outcome. This finding limits the interpretation of a passing score.\par
\section{Related work}
Long-term conversational memory benchmarks examine information across extended interactions. LoCoMo includes question answering, event summarization, and multimodal dialogue generation~\citep{locomo}. LongMemEval examines extraction, multi-session reasoning, temporal reasoning, knowledge updates, and abstention~\citep{longmem}. These constructs overlap with remembering changing preferences. Memory Conflict instead emphasizes the resulting action, its side effects, and who receives private information. Its short synthetic tasks do not replace those broader memory evaluations.\par
Tool-agent evaluation offers a complementary perspective. The $\tau$-bench design checks the database state after an interaction against an annotated goal~\citep{tau}. Memory Conflict also checks state after actions, but focuses on household preference conflicts. Its question tool records delivery rather than a user's reply. It therefore lacks the interactive confirmation process needed to establish real consent. This difference matters when interpreting authority scores.\par
PayBench separates unsafe action from unnecessary refusal through matched commercial scenarios~\citep{paybench}. We use its construct-first presentation and attention to control layers as a design reference. Memory Conflict does not contain matched safe and unsafe lookalikes. All evaluated tasks require a final booking. A refusal baseline can reveal failure to act, but cannot estimate false refusal on a separate benign population.\par
BetterBench treats benchmark quality as a lifecycle problem~\citep{betterbench}. The Agentic Benchmark Checklist highlights how task setup and reward rules can distort measured performance~\citep{abc}. Our study follows that concern by testing the score with deliberately simple behavior. A script with access to the answer serves as a ceiling check. An empty script serves as a floor check. Neither supplies a competitive model baseline.\par
A prospective plan helps separate planned comparisons from changes made after observation~\citep{prereg}. We record the conditions, exclusions, and simulated grid before collecting outcomes. This record is a local versioned plan, not an independently timestamped registration. The distinction prevents a stronger claim of preregistration than the evidence supports.\par
\section{Experiment design}
\subsection{Task and observation model}
Each episode describes a fictional household, its members, stored memories, service options, existing bookings, and a requester. Memories include an owner and a private or shared scope. Some carry an expiry date or a link to the memory they replace. The assistant can read memory, search options, inspect bookings, ask a member, book, cancel, and write memory. The environment exposes member authority and the current date.\par
A memory read includes an active flag. A record becomes inactive when its expiry precedes the current date or an active record supersedes it. This feature reduces the burden of temporal inference. Success on these tasks cannot establish that an assistant can infer validity from unstructured conversation alone. The experiment concerns the use of validity information once it is available.\par
The environment allows booking without a prior question. It also allows information to be written to shared memory. The grader evaluates these actions after replay. The simulator thus exposes mistakes instead of preventing every mistake at execution time. A deployed tool with hard enforcement would present a different task. We do not treat post hoc scoring as proof of an enforcement boundary.\par
The expected outcome specifies the correct option, bookings to cancel, whether clarification is required, and phrases that indicate a privacy breach. It also lists options associated with stale information. The scripted reference uses these expectations. A learned agent should not receive them. Access to this information is deliberate in our instrument checks and is incompatible with a claim of autonomous problem solving.\par
\subsection{Scripted conditions and controls}
The primary study uses \ConditionCount{} scripted conditions on the same \TaskCount{} tasks. Empty, constant, echo, and refusal conditions emit text without booking. The first-tool condition calls the first available tool once and then stops. It tests whether an observation alone earns credit. These controls require no model inference and make no claims about language competence.\par
The packaged reference follows supplied actions. The packaged bad policy chooses an unsuitable option or omits a required action, depending on the task. We then modify reference traces to isolate several failures. No-ask removes all questions. No-cancel removes cancellation calls. Stale-choice replaces a booking with a listed stale option where one exists. Leak copies private memory into shared memory. Unaffected tasks retain their reference actions.\par
Late-ask moves every question after the other reference actions. It preserves the question recipient, text, and final booking state. This is a timing probe, not a repair. A score that measures approval before action should distinguish it from the reference. A score that checks only final state and question presence may not. We report the observed result without redefining the metric after the run.\par
Every condition runs twice with the same runner seed. The runner allows at most twelve turns and twelve tool calls. Retries are disabled, execution is sequential, and the inference budget is zero. Repeated reports retain outputs and tool events. Semantic comparison omits wall-clock timing because timing is not part of the task outcome. Repeated success does not create additional independent tasks.\par
\subsection{Planned analyses and stopping rule}
The first hypothesis is that reference actions score above empty output. The second is that deleting required questions or cancellations fails the corresponding tasks. The simulated hypothesis is that perfect repair improves pass rates as injected errors increase. These are directional instrument checks. We do not perform a search for the best prompt or a significance-based stopping procedure.\par
We retain empty outputs and refusals as failures. We exclude no task based on its score. An infrastructure failure stops the study instead of becoming a model failure. The study ends after the fixed grid and repeated scripted conditions. Free hosted inference is not configured for this study, and provider availability is not tested. No paid inference, human judgment, or learned-model score enters the analysis.\par
\section{Coverage}
The public set contains \TaskCount{} episodes in \TrapCount{} conflict types. Figure~\ref{fig:coverage} shows the counts. These are designed cases from a seeded synthetic generator, not sampled household histories. Variations change members, options, dates, and wording within a limited set of patterns. A large number of replayed traces therefore does not imply broad coverage of household behavior.\par
\begin{figure}[ht]\centering\begin{tikzpicture}\begin{axis}[ybar,ymin=0,ymax=9,ylabel={Episode count},xtick={0,1,2,3,4,5},xticklabels={Authority,Cancel,Expiry,Disagreement,Privacy,Reversal},x tick label style={rotate=22,anchor=east},bar width=22pt,nodes near coords]
\addplot[fill=okblue] table[x=index,y=count,col sep=comma]{data/coverage.generated.csv};
\end{axis}\end{tikzpicture}
\caption{Coverage of the fixed synthetic task set. Privacy has fewer episodes than the other conflict types. Counts describe the design, not estimated household frequencies.}\label{fig:coverage}\end{figure}
The set includes a current request that reverses an older preference, competing household preferences, private information, an expired restriction, an old booking that needs cancellation, and unclear spending authority. Each task has one primary trap label. Some consequences overlap. For example, leaving an old booking active can trigger both a booking failure and the stale-choice flag. Labels and failure indicators are therefore different objects.\par
There is no matched no-conflict control for each task. This prevents a clean estimate of the extra difficulty caused by conflict itself. The paired comparisons in this paper instead hold the task fixed and change the submitted actions. Such comparisons identify scoring sensitivity to those actions. They do not identify the causal effect of confusing memory on a model.\par
The public cases are available before evaluation. Reference actions are also available. Memorization and direct answer lookup can therefore saturate this set. The generator supports other instances, but changing a seed does not create independent conceptual coverage. We evaluate only the fixed public set and make no hidden-test claim. New templates and withheld task families would provide stronger evidence of generalization.\par
\section{Taxonomy of failures}
\paragraph{Stale selection.} An assistant chooses an option linked to an obsolete preference or an expired record. The observed failure is an active booking for a listed stale option. This flag does not reveal which memory the assistant actually read or trusted. The same option can be selected by guessing. We therefore call the measured event a stale choice rather than proof of a mental memory process.
\paragraph{Authority confusion.} An assistant acts for a requester who cannot authorize the booking without asking an adult who can. The current rubric accepts a question to any member with authorization when clarification is required. It does not require a returned answer. It also does not require that the question precede booking. Our results assess this narrower operational condition.
\paragraph{Incomplete replacement.} An assistant creates the right new booking but leaves an obsolete booking active. The final state then contains too many bookings or lacks a required cancellation. Conversely, cancelling an unrelated valid booking also fails. These checks make deletion and addition jointly relevant. Correct selection alone cannot satisfy the full booking condition.
\paragraph{Disclosure across scope.} A private note informs a plan but appears in a message, a final answer, or shared memory visible outside its owner's scope. The grader detects normalized text sequences and configured secret phrases. Copying text into the owner's own private notebook is distinct from sharing it. The relevant property is the recipient or scope, not merely the presence of private text somewhere in the trace.
\paragraph{Wrong choice without stale evidence.} A booked option can be incorrect without belonging to the stale list. Such a trace fails the composite grade but need not set the stale flag. The packaged bad control includes these cases. Reporting only a stale-choice rate would therefore miss other unsuccessful decisions. The composite and diagnostic outcomes should remain visible together.
\paragraph{Procedural mismatch.} The final booking and question set can satisfy the grader despite the wrong sequence of actions. Late-ask is a controlled example. It preserves the score while violating the intended order in the prompt. This is a limit of the measurement rule, not evidence that the changed procedure is safe. It motivates a future trace-order requirement.
\section{Metrics and scoring}
\subsection{Composite success}
Let $B$ denote the correct active booking state, $C$ all required cancellations, and $W$ the absence of wrongful cancellation. Let $L$ denote a detected leak, $S$ a listed stale option remaining active, and $A$ the question condition. The native binary grade is\par
\begin{equation}Y=\mathbf{1}\{B\land C\land W\land\neg L\land\neg S\land A\}.\end{equation}
When an option is required, $B$ means that exactly one active booking remains and its option matches the expected option. The grade does not award partial credit for a correct explanation. It also does not treat a safe but unhelpful refusal as completion. This strict conjunction is appropriate for a task that requires all specified conditions, but hides the source of failure unless components are reported.\par
The diagnostic stale and leak rates use all tasks as the denominator. Question success uses only tasks that require clarification. Using all tasks for question success would include cases where the requirement is vacuously true. The conditional denominator better answers whether the assistant asks when necessary. A question success rate is still not a consent rate because the environment records no affirmative reply.\par
The privacy check lowercases text, removes punctuation, and applies limited word normalization. It checks short sequences from private memories and configured phrases against outbound content. This detects the deliberate copying in our leak control. It does not establish semantic confidentiality. A paraphrase may evade the matcher, while an innocent overlap may trigger it. No human-labeled precision or recall estimate is available.\par
\subsection{Uncertainty and dependence}
For a binary task outcome with $k$ successes among $n$ tasks, we report the Wilson interval at the nominal 95\% level. This interval avoids a zero-width interval at an observed boundary. It describes uncertainty under a hypothetical exchangeable-task sampling model. Our actual task set is fixed and template-based. The interval therefore is not a guarantee about households, service domains, or future templates.\par
The repeated run checks reproducibility rather than statistical sampling. We report it separately and never double the task denominator. For a targeted mutation, unaffected tasks can dominate the aggregate score. We therefore include stratum results. A high aggregate score for no-ask cannot establish authority handling when the affected authority stratum fails.\par
The simulated study uses a different sampling unit: a seeded simulated agent evaluated on the entire fixed task set. We first compute each agent's mean score. A percentile bootstrap resamples those agent means with replacement. The paired control contrast resamples within-agent differences. It preserves the common error draws used across control conditions and does not treat each replay as a new household.\par
\section{Results}
\subsection{Floor, reference, and targeted omissions}
Table~\ref{tab:results} reports every scripted condition. The reference achieves \Good{}. Empty output achieves \Empty{}. Constant output, prompt echo, refusal, and a single observation also fail to complete the booking task. These results support the planned floor-versus-reference check. They do not demonstrate robust natural-language understanding because the reference uses privileged actions.\par
\begin{table}[ht]\centering\begin{tabular}{lcc}\toprule
Script & Pass rate [95\% CI] & Passed / tasks\\\midrule
empty & 0.000 [0.000, 0.088] & 0/40\\
constant & 0.000 [0.000, 0.088] & 0/40\\
echo & 0.000 [0.000, 0.088] & 0/40\\
refuse & 0.000 [0.000, 0.088] & 0/40\\
first-tool & 0.000 [0.000, 0.088] & 0/40\\
packaged-good & 1.000 [0.912, 1.000] & 40/40\\
packaged-bad & 0.000 [0.000, 0.088] & 0/40\\
no-ask & 0.825 [0.681, 0.913] & 33/40\\
no-cancel & 0.825 [0.681, 0.913] & 33/40\\
stale-choice & 0.300 [0.181, 0.454] & 12/40\\
leak & 0.875 [0.739, 0.945] & 35/40\\
late-ask & 1.000 [0.912, 1.000] & 40/40\\\bottomrule\end{tabular}
\caption{Native composite success for all scripted controls. Brackets contain 95\% Wilson intervals over the fixed task set.}\label{tab:results}\end{table}
Removing questions yields overall pass rate \NoAsk{}. On tasks that require a question, clarification success becomes \AskFail{}. Removing cancellations yields overall pass rate \NoCancel{}. The corresponding stratum results appear in the appendix. The similarity of the two aggregate scores reflects the designed coverage. It does not imply equivalent harms or equivalent model difficulty.\par
Replacing choices with listed stale options yields \Stale{}. Copying private memories into shared memory yields \Leak{}. The latter aggregate remains high because most tasks have no private note to copy. This illustrates why aggregate completion alone can conceal a complete failure on a small safety-relevant stratum. The component table makes the detected disclosures explicit.\par
\begin{figure}[ht]\centering\begin{tikzpicture}\begin{axis}[ymin=0,ymax=1.1,ylabel={Composite pass rate},xtick={0,1,2,3,4,5,6,7,8,9,10,11},xticklabels={Empty,Constant,Echo,Refuse,First tool,Reference,Bad,No ask,No cancel,Stale,Leak,Late ask},x tick label style={rotate=45,anchor=east},height=6cm]
\addplot[only marks,mark=*,okblue,error bars/.cd,y dir=both,y explicit] table[x=index,y=value,y error minus=minus,y error plus=plus,col sep=comma]{data/policies.generated.csv};
\end{axis}\end{tikzpicture}
\caption{Scripted pass rates with 95\% Wilson intervals. Reference and late-ask scores coincide, while omissions fail affected tasks. These points are instrument controls, not learned-model results.}\label{fig:policies}\end{figure}
\subsection{Timing null and its meaning}
Late-ask achieves \Late{}, the same composite result as the reference. Conditional question success is \AskPass{}. This is a null difference under the current grade. The changed ordering does not change the scored final state or the set of members asked. We retain the null rather than interpreting it as successful authority preservation.\par
The observed invariance supports a specific conclusion: the current metric does not distinguish asking before booking from asking afterward in these traces. It does not support a broader claim that all temporal failures are invisible. For example, a different sequence can change the final booking state. The boundary is the information retained by the scoring predicate, not time in general.\par
No affirmative response is needed for either trace. The tool acknowledges that the message was delivered, and the grader checks the recipient's authority. A real approval flow would need a reply, its relation to the proposed action, and a check that it preceded execution. Those stages are outside this benchmark version. A perfect score should therefore be described as rubric compliance, not verified consent.\par
\subsection{Reproducibility}
The scripted study completes \EpisodeCount{} runner episodes, including the repeated conditions. The semantic traces and scores agree across repeats. This establishes stability for the fixed scripted inputs and frozen environment. It does not establish reliability under stochastic model responses, changing tools, or unavailable services. No inference cost is incurred by the scripts.\par
\section{Extra virtual experiment: simulated control ablation}
\subsection{Simulated population and intervention}
The additional experiment is explicitly simulated. Each of \SeedCount{} seeds defines an artificial agent's error pattern across the fixed tasks. We draw a reproducible uniform value per seed and task. An error is injected when that value falls below a specified probability. We use probabilities from zero to one in quarter steps. Common draws preserve pairing across probabilities and controls.\par
For authority tasks, the simulated error removes the question. For cancellation tasks, it removes cancellation. For private-note tasks, it copies the private text into shared memory. The remaining trap types receive a listed stale choice. This mapping selects a known failing action for each primary trap. It is an assumed error model, not a distribution fitted to observed agents.\par
We compare no repair, perfect authority repair, perfect privacy repair, perfect cancellation repair, and all repairs. A targeted repair restores the reference trace only on its target trap. All repairs also restore stale choices. The controls know the expected answer and make no false-positive interventions. They are oracle upper bounds for the selected failure pattern. They are not implementable defenses evaluated on realistic input.\par
Every simulated trace is scored by the same public grading path. The sweep contains \SimCount{} such traces. The saved records identify the seed, task, probability, control, selected trace policy, score, and diagnostic flags. The trace policy and original task fully determine the replay input. This permits reconstruction without storing identical reference actions repeatedly.\par
\subsection{Simulated findings}
At injected probability $q=0.5$, simulated no-control success is \SimNone{}. Simulated all-repair success is \SimAll{}. The paired simulated gain is \SimGain{}. These intervals resample simulated agent seeds, not tasks or real users. The all-repair interval collapses because every repaired trace is the reference by construction. It does not mean a real control has no uncertainty.\par
\begin{figure}[ht]\centering\begin{tikzpicture}\begin{axis}[xlabel={Injected error probability (simulated)},ylabel={Pass rate (simulated)},ymin=0,ymax=1.05,legend style={font=\footnotesize,at={(.5,-.28)},anchor=north,legend columns=3},height=6cm]
\foreach \control/\label in {none/None,authority/Authority,privacy/Privacy,cancellation/Cancellation,all/All repairs}{\addplot+[error bars/.cd,y dir=both,y explicit] table[x=q,y=value,y error minus=minus,y error plus=plus,col sep=comma]{data/sim-\control.generated.csv};\expandafter\addlegendentry\expandafter{\label}}
\end{axis}\end{tikzpicture}
\caption{Simulated control ablation with 95\% intervals from resampling simulated agent means. Perfect targeted repairs recover only their covered trap; all repairs restore the reference by construction.}\label{fig:simulation}\end{figure}
The simulated curves separate as the injected error probability rises. A targeted control can repair only the share of tasks assigned to its trap. Privacy repair has less potential coverage than a control assigned to a larger stratum. This is a property of the task mixture. It is not a finding that privacy matters less or that privacy failures are easier to prevent.\par
At zero injected error, the simulated controls have no fault to repair. The paired difference is null. At full injected error, the simulated all-repair condition restores the reference and the no-control condition retains the chosen failures. The endpoints follow directly from the construction. The intermediate points verify that the execution and aggregation preserve the planned pairing.\par
\subsection{Analytic check and scope}
Let $w_j$ be the fraction of fixed tasks covered by a perfect repair $j$. Under the simulated model, all injected mutations fail and all reference traces pass. The expected no-control score is $1-q$. A single perfect repair has expected score $1-q+qw_j$. This relation explains why a curve can improve without any change in memory reasoning. The measured finite-seed curves need not equal these expectations exactly.\par
The simulation supplies a sensitivity analysis of the score and task mixture. It does not estimate field prevalence, control latency, user burden, or the cost of unnecessary intervention. Errors are assigned independently of language complexity. The same task set is reused for every simulated agent. Future model studies must replace these assumptions with observed behavior before recommending a control.\par
\section{Roadmap}
The first priority is to separate question delivery from authorization. A revised environment should expose a reply state and tie approval to the requested booking. The grader should check the ordering of request, approval, and execution. A consent check should also reject a reply that applies to another option. The present timing probe supplies a concrete regression case for that extension.\par
The second priority is matched controls. Each conflict case should have a version where acting directly is appropriate. Such pairs would permit separate estimates of unsafe action and unnecessary interruption. The pair should preserve surface wording and option difficulty where possible. Otherwise, a difference could reflect task complexity rather than the conflict rule.\par
The third priority is broader coverage with explicit sampling claims. Combined traps would test whether privacy and authority constraints interact. Alternative domains would test whether the same action rules transfer beyond the existing services. New household structures would require careful review of authority assumptions. A larger seed sweep alone cannot replace these design changes.\par
The fourth priority is external validation. Human reviewers should inspect whether the task rubric captures a defensible action. Independent annotations of privacy disclosures would help estimate false positives and missed paraphrases. Learned models should then run without expected outcomes or reference traces. Their raw actions should remain available for error analysis, including refusals and malformed tool calls.\par
\section{Caveats and responsible interpretation}
\paragraph{Construct coverage.} The benchmark models explicit local constraints, not the full meaning of household preference or consent. Its authority labels are assigned facts. They should not be used to infer authority in a real household. The expired-allergy cases are synthetic date tests, not advice to disregard an actual medical restriction.
\paragraph{Score validity.} Final-state checks miss the demonstrated timing distinction. Privacy matching is lexical and has no independently measured error rate. A listed stale option is an observable choice, not proof of the assistant's internal reasoning. These limitations constrain the labels used for all reported components.
\paragraph{Sampling and uncertainty.} The public set is small and generated from related templates. Wilson intervals do not account for an unknown population of household situations. Simulated bootstrap intervals capture variation in selected artificial seeds only. Adjacent seeds of a fixed pseudorandom generator are a reproducibility device, not sampled human identities. Boundary intervals can collapse in the simulated experiment.
\paragraph{Baselines and access.} The scripted reference and repairs use privileged task information. They validate the scoring path and set an artificial ceiling. No learned model, human participant, or provider comparison is present. Free hosted inference was not configured; we make no claim that free models were unavailable elsewhere. The scripts require no model credentials.
\paragraph{Public exposure.} Public tasks and reference actions permit direct memorization. An unpublished seed would hide particular instances but not the generator's concepts. Claims about generalization require withheld templates or independently authored cases. The present result is a dated audit of a fixed benchmark definition.
\paragraph{Study governance.} The prospective plan was committed before data collection, but has no independent deposit. The author maintains the benchmark and selects the controls. This is an internal audit rather than an independent certification. No human study, demographic comparison, or real booking occurs. The paper and figures use CC-BY-4.0; the benchmark retains its existing software and task license.
\section{Citation and availability}
Code, tasks, raw scripted reports, simulated records, and reproduction instructions are available at \url{https://github.com/pooriaarab/benchmark-memory-conflict}. Cite this paper by its title and author, and identify the evaluated benchmark revision when comparing results. No archival publication or venue acceptance is claimed.\par
\section{Authors}
Pooria Arab is an independent researcher and the benchmark maintainer. AI assistance supported implementation and drafting. The author is responsible for the scientific claims and release decisions. No external human review or independent replication is claimed in this packet.\par
\bibliographystyle{plainnat}
\bibliography{references}
\appendix
\section{Detailed diagnostic outcomes}
Table~\ref{tab:details} separates stale choice, detected disclosure, and required-question success. The first two columns use all tasks. The last column uses only required-question tasks. The different denominators are intentional. For example, a trace can avoid a disclosure by doing nothing while still failing the booking. A low diagnostic failure rate alone does not establish useful task completion.\par
\begin{table}[ht]\centering\resizebox{\linewidth}{!}{\begin{tabular}{lccc}\toprule
Script & Stale choice & Privacy leak & Required question\\\midrule
packaged-good & 0.000 [0.000, 0.088] & 0.000 [0.000, 0.088] & 1.000 [0.646, 1.000]\\
packaged-bad & 0.300 [0.181, 0.454] & 0.125 [0.055, 0.261] & 0.000 [0.000, 0.354]\\
no-ask & 0.000 [0.000, 0.088] & 0.000 [0.000, 0.088] & 0.000 [0.000, 0.354]\\
no-cancel & 0.175 [0.087, 0.319] & 0.000 [0.000, 0.088] & 1.000 [0.646, 1.000]\\
stale-choice & 0.700 [0.546, 0.819] & 0.000 [0.000, 0.088] & 1.000 [0.646, 1.000]\\
leak & 0.000 [0.000, 0.088] & 0.125 [0.055, 0.261] & 1.000 [0.646, 1.000]\\
late-ask & 0.000 [0.000, 0.088] & 0.000 [0.000, 0.088] & 1.000 [0.646, 1.000]\\\bottomrule\end{tabular}
}\caption{Diagnostic rates with 95\% Wilson intervals. Question success means a qualified member was asked; it does not mean approval was received.}\label{tab:details}\end{table}
The no-cancel control can set the stale flag even though the new booking is correct. The obsolete booking remains active. This shows that diagnostic flags need not be disjoint. Adding their rates would double-count some episodes and omit other booking errors. The composite is a conjunction of requirements, not the complement of a sum of independent failure probabilities.\par
The packaged bad policy is a deliberately flawed program with expected-outcome access. Its mixture of wrong choices and missing questions differs from a learned model's possible failures. It is useful as a rejection check. It is not a meaningful estimate of how often an untrained or careless assistant would disclose information.\par
\section{Stratum-level scoring sensitivity}
Table~\ref{tab:strata} gives task-level pass rates within each primary trap. This view separates the breadth of a mutation from its effect where it applies. No-ask is unchanged on tasks without a question requirement. No-cancel is unchanged on tasks without an old booking to remove. The leak mutation is unchanged when there is no private memory.\par
\begin{table}[ht]\centering\resizebox{\linewidth}{!}{\begin{tabular}{lcccc}\toprule
Conflict type & No ask & No cancel & Stale choice & Leak\\\midrule
authority unclear & 0.000 [0.000, 0.354] & 1.000 [0.646, 1.000] & 1.000 [0.646, 1.000] & 1.000 [0.646, 1.000]\\
cancel old booking & 1.000 [0.646, 1.000] & 0.000 [0.000, 0.354] & 0.000 [0.000, 0.354] & 1.000 [0.646, 1.000]\\
expired allergy & 1.000 [0.646, 1.000] & 1.000 [0.646, 1.000] & 0.000 [0.000, 0.354] & 1.000 [0.646, 1.000]\\
household disagreement & 1.000 [0.646, 1.000] & 1.000 [0.646, 1.000] & 0.000 [0.000, 0.354] & 1.000 [0.646, 1.000]\\
private note & 1.000 [0.566, 1.000] & 1.000 [0.566, 1.000] & 1.000 [0.566, 1.000] & 0.000 [0.000, 0.434]\\
reversed preference & 1.000 [0.646, 1.000] & 1.000 [0.646, 1.000] & 0.000 [0.000, 0.354] & 1.000 [0.646, 1.000]\\\bottomrule\end{tabular}
}\caption{Composite pass rates by primary trap, with 95\% Wilson intervals. The wide intervals reflect small strata.}\label{tab:strata}\end{table}
A stale substitution can apply to more than the trap named reversal. The expected outcome also lists stale options for other contexts. The mutation follows that explicit list rather than guessing from a label. This choice makes the audit reproducible and explains why its aggregate failure count differs from a single-stratum ablation.\par
These intervals should not be used to rank trap difficulty. The controls are designed to satisfy or violate known predicates, and the small strata share templates. Difficulty is a relation between a task distribution and a solver. This experiment has no learned solver and no representative distribution of households.\par
\section{Paired simulated contrasts}
\begin{table}[ht]\centering\begin{tabular}{lc}\toprule
Injected probability & All minus none [95\% CI]\\\midrule
0.00 & 0.000 [0.000, 0.000]\\
0.25 & 0.269 [0.246, 0.293]\\
0.50 & 0.522 [0.498, 0.545]\\
0.75 & 0.778 [0.755, 0.804]\\
1.00 & 1.000 [1.000, 1.000]\\\bottomrule\end{tabular}
\caption{Simulated all-repair minus no-control differences. The 95\% percentile intervals resample paired simulated-agent means.}\label{tab:simulation}\end{table}
The simulated contrast holds the seed and task fixed before comparing controls. If an error draw falls above the selected probability, neither condition needs a repair. If it falls below, all-repair restores the reference. This pairing reduces variation caused solely by different artificial error assignments. It does not address uncertainty about whether the chosen errors represent real agent behavior.\par
The simulated bootstrap draws whole agent means with replacement. It uses a fixed resampling seed and a fixed number of draws. Repeating the experiment reproduces the same semantic records and intervals. Different simulation generators could change finite-seed estimates. The analytic relation in the main text remains the more direct description of the assumed mechanism.\par
\section{Worked action distinctions}
For an authority case, compare a reference trace that reads memory, asks an adult, and books with one that reads memory, books, and asks. Both end with the same option and the same member in the asked set. The current score therefore cannot separate them. To measure prior authorization, the state would need to retain the approval event and its position relative to booking.\par
For a replacement case, compare cancelling an obsolete booking before a new booking with leaving the obsolete booking active. The second trace has an extra active booking and lacks a required cancellation. The current score distinguishes these final states. This example explains why final-state scoring can correctly detect one procedural omission while missing another temporal distinction.\par
For a privacy case, compare keeping a note within its owner's private scope with copying that note into shared memory. The latter exposes matching text outside the owner's scope. The leak control uses direct copying so its intended effect is unambiguous to the lexical matcher. A semantic paraphrase control would test a different limitation and requires a separate plan and annotated expectations.\par
\section{Evidence boundaries for future studies}
An independent replication should first reproduce the scripted controls. It should then verify that learned agents cannot access expected outcomes. The model identifier, provider, settings, task revision, and raw tool events should accompany any new scores. A timeout should remain distinct from a valid refusal. A repeated run should not silently replace a failed earlier attempt.\par
A control evaluation should report both prevented mistakes and unnecessary interventions. Perfect restoration is only an upper-bound instrument. A practical control can block a valid request, miss a paraphrase, or ask the wrong person. The next study should measure those costs on matched tasks before describing a safety-autonomy tradeoff. The current simulated curves cannot supply that evidence.\par
\end{document}
